%%%%%%%%%%%%%%%%%%%%%%%%%%%%%%%%%
%
%       EXERCÍCIO
%
%%%%%%%%%%%%%%%%%%%%%%%%%%%%%%%%%

\ifdefstring{\atividade}{prova}{%
    \ifdefstring{\modo}{objetivo}{%
        \renewcommand{\valorquestao}{\ValorQObj\ ponto}
    }{%
        \renewcommand{\valorquestao}{\ValorQDisc\ pontos}
    }%
}%

\begin{exercicioBanco}[\valorquestao]
Considere a função
%
\[
    \begin{array}{rcl}
        f: \; \bbR & \to & \bbR \\
        x & \mapsto & \left\{
            \begin{array}{cl}
                2+x, & \text{ se }\; x \le -1, \\[2pt]
                1, & \text{ se }\; -1 < x < 1, \\[2pt]
                -2x+3, & \text{ se }\; x \ge 1,
            \end{array}
        \right.
    \end{array}
\] 
%
Calcule \(f(-1) - f(0) + f\left(\dfrac{3}{2}\right)\).

% Define as alternativas
\newcommand{\alternativas}{%
    \begin{center}
        \begin{tabularx}{\textwidth}{XXXXX}
            (a) \(-1\). &
            (b) \(1\). &
            (c) \(0\). &
            (d) \(-2\). &
            (e) \(2\).
        \end{tabularx}
    \end{center}
}

% Define a resposta correta
\newcommand{\resposta}{C}

% Lógica condicional para exibição
\ifdefstring{\atividade}{lista}{%
    \alternativas
    \vspace{0.5em}
    
    \noindent\textbf{Resposta:} letra \textbf{\resposta}.
}{%
    \ifdefstring{\modo}{objetivo}{%
        \alternativas
    }{%
        % modo = discursiva → não mostra alternativas
    }
}
\end{exercicioBanco}

